\uuid{nwGM}
\titre{Base et dimension}
\auteur{hanani ; mzari}
\organisation{exo7}
\licence{CC BY-NC-SA 4.0}
\datecreate{2020-01-20}
\niveau{L1}
\module{Algèbre}
\chapitre{Espace vectoriel}
\sousChapitre{Base}
\difficulte{3}
\theme{Base et dimension}
\exo7qcmid{lille-2.espaces-vectoriels.045}

\contenu{
\begin{question}
\qclassification{106.04}
\qtags{motcle=Bases de $\Rr$}
Quelles sont les assertions vraies ?
\begin{answers}
\good{$\{1\}$ est une base de $\Rr$ comme $\Rr$-espace vectoriel}.
\good{$\{\sqrt 2\}$ est une base de $\Rr$ comme $\Rr$-espace vectoriel}.
\bad{$\{1,\sqrt 2\}$ est une base de $\Rr$ comme $\Rr$-espace vectoriel}.
\bad{$\{1, \sqrt 2\}$ est une base de $\Rr$ comme $\Qq$-espace vectoriel}.
\end{answers}
\begin{explanations} $\Rr$ est un  $\Rr$-espace vectoriel de dimension 1, donc pour tout $\alpha \in \Rr^*$, $\{\alpha\}$
est une base de $\Rr$.
\vskip2mm
$\{1, \sqrt 2\}$ n'est pas une base de $\Rr$ comme  $\Qq$-espace vectoriel. En effet, sinon, il existe $\alpha, \beta \in \Qq$ tels que $\sqrt 3= \alpha+ \beta \sqrt 2$. En considérant le carré de cette égalité, on déduit que $\sqrt 2$ est un rationnel, ce qui est absurde.
\end{explanations}
\end{question}
}
